
\needspace{3\baselineskip}
\color{uniblue}{\subsection{Коммутационные аппараты (КА)}}
\color{black}

\begin{enumerate}[label=\arabic{section}.\arabic{subsection}.\arabic{enumi}, labelsep=4pt, leftmargin=0pt, itemindent=57pt, itemsep=0pt, parsep=5pt] % если что можно крутить itemsep

%%%%%%%%%%%%%%%%%%%%%%%%%%%%%%%%%%%%%%%%%%%%%%%%%%%%%%%%%%%%%%%%%%%%% ОБЩИЕ СВЕДЕНИЯ %%%%%%%%%%%%%%%%%%%%%%%%%%%%%%%%%%%%%%%%%%%%%%%%%%%%%%%%%%%%%%

\item Общие сведения

\begin{enumerate}[label=\arabic{section}.\arabic{subsection}.\arabic{enumi}.\arabic*, labelsep=4pt, leftmargin=0em, itemindent=65pt, parsep=0pt]

\item
Модуль коммутационных аппаратов предназначен для реализации представлений (моделей) высоковольтных коммутационных аппаратов (заземляющий нож, разъединитель, выключатель в составе защищаемого присоединения) в устройстве.

\item
В устройстве модуль коммутационных аппаратов представлен функциональным блоком «КА», который состоит из одной модели силового выключателя -- функция «Выключатель».

\item
Ввод функционального блока «КА» в работу на этапе параметрирования устройства осуществляется программным переключателем SGF1 «Ввод\_функции». Введенное состояние функционального блока соответствует активному состоянию сигнала «КА: Ввод».

\item
Функциональный блок «КА» может быть выведен из работы оперативно путем активации сигнала «ОВ КА», о чем сигнализирует активное состояние сигнала «КА: Оперативный вывод».

\item
Логика управления режимами работы функционального блока «КА» выполнена в соответствии с рисунком \ref{ka:img_gen}. В таблице \ref{ka:tbl_gen} приведены параметры, необходимые для настройки функционального блока «КА».

\vspace{2mm}
\begin{figure}[H]
\centering

\includegraphics[width=0.9\textwidth,height=0.9\textheight,keepaspectratio]{\fbpath/46. КА В/_latex/img/1.pdf}
\captionof{figure}{Логика управления режимами функционального блока «КА»}\label{ka:img_gen}
\end{figure}

\small
\begin{longtable}{|>{\centering\arraybackslash}m{5.3cm}|>{\centering\arraybackslash}m{3.3cm}|>{\centering\arraybackslash}m{4.2cm}|>{\centering\arraybackslash}m{1.8cm}|>{\centering\arraybackslash}m{1cm}|}
\caption{Параметры для настройки функционального блока «КА»\hfill\vspace{-0.5\baselineskip}}\label{ka:tbl_gen}\\ 
\hline
\rowcolor{gray!30}
Параметр (Параметр на ИЧМ) & Условное обозначение на схеме & Значение/ Диапазон & Единица измерения & Шаг \\ 
\hline
\endfirsthead
\caption*{\hspace{3pt}\emph{Продолжение таблицы \ref{ka:tbl_gen}\hfill\vspace{-0.5\baselineskip}}} \\ % сделано по ГОСТ 2.105 п.6.8.7
\hline
\rowcolor{gray!30}
Параметр (Параметр на ИЧМ) & Условное обозначение на схеме & Значение/ Диапазон & Единица измерения & Шаг \\ 
%\hline
\endhead
%\hline
\endfoot
%\hline
\endlastfoot
%==+t1*LLN0|T_SwitchDevice
\centering Ввод функции в работу (Ввод\_функции) & \centering SGF1 & \centering 0 = Не предусмотрено\\1 = Предусмотрено & \centering -- & \centering \arraybackslash -- \\
%===t1
\end{longtable}
\normalsize

\end{enumerate}
\FloatBarrier % Все float'ы ДО этой строки будут размещены здесь.
%%%%%%%%%%%%%%%%%%%%%%%%%%%%%%%%%%%%%%%%%%%%%%%%%%%%%%%%%%% ВЫКЛЮЧАТЕЛЬ %%%%%%%%%%%%%%%%%%%%%%%%%%%%%%%%%%%%%%%%%%%%%%%%%%%%%%%%%%

\item Выключатель (В)

\begin{enumerate}[label=\arabic{section}.\arabic{subsection}.\arabic{enumi}.\arabic*, labelsep=4pt, leftmargin=0em, itemindent=65pt, parsep=0pt]

\item
В функциональном блоке «КА» силовой выключатель представлен функцией «Выключатель».

\item
При приеме входных сигналов от выходной логики отключения выключателя стороны ВН трансформатора и функции УРОВ («ЛО ВН / ЛО: Отключить аварийно», «УРОВ: Сраб на себя») или от логики оперативного управления («КП / УВ: Отключить», «Опер. отключить В» -- от внешнего КП) формируется команда «В: Отключить (реле)» на внутреннее выходное реле устройства с целью коммутации выходной цепи отключения ЭМО выключателя. При этом возможна выдача как импульсного сигнала (режим «Импульсный» программного переключателя SGF2 «Режим\_откл», время импульса определяется уставкой Т2 «Тимп\_откл»), так и длительного сигнала (режим «Длительный» программного переключателя SGF2 «Режим\_откл»). 
В режиме выдачи импульсного сигнала выдаваемый импульс (а также и выдаваемый сигнал в длительном режиме) ограничивается активизацией сигнала отключенного состояния выключателя «В: Отключен» с учетом или без учета сигналов «Работа ЭМО1» и «Работа ЭМО2» (в зависимости от положения программного переключателя SGF3 «Контр\_раб\_ЭМО»).
\item
При приеме входного сигнала от логики оперативного управления («КП / УВ: Включить», «УВ / УВ: Включить») формируется команда «В: Включить (реле)» на внутреннее выходное реле устройства с целью коммутации выходной цепи включения ЭМВ выключателя. При этом возможна выдача как импульсного сигнала (режим «Импульсный» программного переключателя SGF4 «Режим\_вкл», время импульса определяется уставкой Т3 «Тимп\_вкл»), так и длительного сигнала (режим «Длительный» программного переключателя SGF4 «Режим\_вкл»).
В режиме выдачи импульсного сигнала выдаваемый импульс (а также и выдаваемый сигнал в длительном режиме) ограничивается активизацией сигнала включенного состояния выключателя «В: Включен» с учетом или без учета сигнала «Работа ЭМВ» (в зависимости от положения программного переключателя SGF5 «Контр\_раб\_ЭМВ»). При этом возможна задержка активации сигнала «В: Включен» (определяется уставкой Т4 «Тпрод\_вкл») для сброса импульсного или длительного сигнала на включение (для исключения ситуаций так называемого «опрокидывания» выключателя при раннем съёме команды «В: Включить (реле)» -- это характерно для некоторых видов масляных выключателей).
\item
Для снятия команд «Включить (реле)» или «Отключить (реле)» при отказе высоковольтного выключателя необходимо принудительно обесточить цепи управления и подать команду «Сброс».
\item
При выводе функции из работы (отсутствие выходного сигнала «В: Ввод») происходит сброс команд включения/отключения выключателя.
\item
В логике функции реализована блокировка на прохождение команды включения выключателя «В: Включить (реле)» при приеме сигнала аварийного отключения от защит.
\item
Логика функции формирует сигналы состояния выключателя «В: Отключен», «В: Включен», «В: Промеж. положение» и «В: Неиспр. положение», при этом сигнал «В: Неиспр. положение» выдается с выдержкой времени, задаваемой уставкой Т1 «Тнеиспр\_В».

\item
Ввод функции  в работу на этапе параметрирования устройства осуществляется программным переключателем SGF1 «Ввод\_функции». Введенное состояние функции соответствует активному состоянию сигнала «В: Ввод».

\item
Алгоритм работы функции выполнен в соответствии с рисунком \ref{ka:img_v}. В таблице \ref{ka:tbl_v} приведены параметры, необходимые для настройки функции.

\vspace{2mm}
\begin{figure}[H]
\centering
\includegraphics[width=0.8\textwidth,height=0.8\textheight,keepaspectratio]{\fbpath/46. КА В/_latex/img/v.pdf}
\captionof{figure}{Функциональная схема логики функции «Выключатель»}\label{ka:img_v}
\end{figure}

\small
\begin{longtable}{|>{\centering\arraybackslash}m{5.3cm}|>{\centering\arraybackslash}m{3.3cm}|>{\centering\arraybackslash}m{4.2cm}|>{\centering\arraybackslash}m{1.8cm}|>{\centering\arraybackslash}m{1cm}|}
\caption{Параметры для настройки функции «Выключатель»\hfill\vspace{-0.5\baselineskip}}\label{ka:tbl_v}\\ 
\hline
\rowcolor{gray!30}
Параметр (Параметр на ИЧМ) & Условное обозначение на схеме & Значение/ Диапазон & Единица измерения & Шаг \\ 
\hline
\endfirsthead
\caption*{\hspace{3pt}\emph{Продолжение таблицы \ref{ka:tbl_v}\hfill\vspace{-0.5\baselineskip}}} \\ % сделано по ГОСТ 2.105 п.6.8.7
\hline
\rowcolor{gray!30}
Параметр (Параметр на ИЧМ) & Условное обозначение на схеме & Значение/ Диапазон & Единица измерения & Шаг \\ 
%\hline
\endhead
%\hline
\endfoot
%\hline
\endlastfoot
%==+t1*XCBR1|T_SwitchDevice	
\centering Ввод функции в работу (Ввод\_функции) & \centering SGF1 & \centering 0 = Не предусмотрено\\1 = Предусмотрено & \centering -- & \centering \arraybackslash -- \\
\hline
\centering Режим отключения В (Режим\_откл) & \centering SGF2 & \centering 0 = Длительный\\1 = Импульсный & \centering -- & \centering \arraybackslash -- \\
\hline
\centering Контроль работы ЭМО (Контр\_раб\_ЭМО) & \centering SGF3 & \centering 0 = Не предусмотрено\\1 = Предусмотрено & \centering -- & \centering \arraybackslash -- \\
\hline
\centering Режим включения В (Режим\_вкл) & \centering SGF4 & \centering 0 = Длительный\\1 = Импульсный & \centering -- & \centering \arraybackslash -- \\
\hline
\centering Контроль работы ЭМВ (Контр\_раб\_ЭМВ) & \centering SGF5 & \centering 0 = Не предусмотрено\\1 = Предусмотрено & \centering -- & \centering \arraybackslash -- \\
\hline
\centering Длительность формирования сигнала о неисправном положении В (Tнеиспр\_В) & \centering T1 & \centering 0,01 ... 10,00 & \centering с & \centering \arraybackslash 0,01 \\
\hline
\centering Длительность формирования команды «Отключить В (реле)» (Тимп\_откл) & \centering T2 & \centering 0,05 ... 10,00 & \centering с & \centering \arraybackslash 0,01 \\
\hline
\centering Длительность формирования команды «Включить В (реле)» (Тимп\_вкл) & \centering T3 & \centering 0,05 ... 10,00 & \centering с & \centering \arraybackslash 0,01 \\
\hline
\centering Время для исключения «опрокидывания» В (Тпрод\_вкл) & \centering T4 & \centering 0,00 ... 1,00 & \centering с & \centering \arraybackslash 0,01 \\
%===t1
\end{longtable}
\normalsize

\end{enumerate}

\end{enumerate}
\FloatBarrier % Все float'ы ДО этой строки будут размещены здесь.